\chapterbackmatter{Solutions to Weinberg Exercises}
\ExerciseEditorialNote

\begin{enumerate}
\WeinbergSolution{1}
Continue the argument of Appendix 24.B with a generator whose kernel contains
at most \(D\) derivatives of the momentum delta function.  Its \(D\)-fold
commutator with translations has the one-particle matrix
\[
 b^{\mu_1\cdots\mu_D}(p)
 =b^{\#\mu_1\cdots\mu_D}
  a^{\nu\mu_1\cdots\mu_D}p_\nu\mathbf 1 ,
\]
where the last \(D\) indices of \(a\) are symmetric.  Preservation of the
mass shell gives
\[
 p_{\mu_1}b^{\mu_1\cdots\mu_D}(p)=0
 \qquad\text{for every }p\text{ with }p^2=0.
\]
Scaling a fixed null vector separates the terms of first and second degree in
\(p\).  A linear form that vanishes on the whole null cone is zero, whereas a
quadratic form that vanishes there is proportional to \(p^2\).  Consequently,
for \(D\geq1\),
\[
 b^{\#\mu_1\cdots\mu_D}=0,\qquad
 a^{(\nu\mu_1)\mu_2\cdots\mu_D}
 =\eta^{\nu\mu_1}c^{\mu_2\cdots\mu_D}.
\]
This is the massless replacement for Eq.~\eqref{eq:24.B.30}.

The possibilities can now be classified by the order \(D\).  For \(D=0\)
one gets translations and generators commuting with all translations; the
latter are internal symmetries.  For \(D=1\),
\[
 a^{\mu\nu}=\omega^{\mu\nu}+\lambda\eta^{\mu\nu},
 \qquad \omega^{\mu\nu}=-\omega^{\nu\mu}.
\]
The antisymmetric part gives the Lorentz generators, while the trace gives a
dilatation \(D\), normalized so that
\([P^\mu,D]=iP^\mu\).

For \(D=2\), symmetry in the last two indices fixes the nonzero top symbol,
up to normalization, to
\[
 a^{\mu\nu\rho}
 =\eta^{\mu\nu}\rho^\rho+\eta^{\mu\rho}\rho^\nu
  -\eta^{\nu\rho}\rho^\mu .
\]
The corresponding four generators are the special conformal generators
\(K^\mu\): taking two commutators of
\([P^\mu,K^\nu]=2i\eta^{\mu\nu}D+2iJ^{\mu\nu}\)
with translations gives precisely this tensor structure.  A symbol of order
\(D>2\) would, under repeated commutation with translations and with its
Lorentz transforms, generate independent symbols of unbounded order unless
its leading tensor vanished.  That contradicts the assumed
finite-dimensional symmetry algebra.  Equivalently, the finite-dimensional
first-order transformations of spacetime that preserve the null cone are
exactly translations, Lorentz transformations, dilatations, and special
conformal transformations.

Finally the Jacobi identities and Lorentz covariance fix all remaining
brackets.  They are Eqs.~\eqref{eq:24.B.34} and
\eqref{eq:24.B.35}.  Internal generators commute with this spacetime algebra
by the same scattering argument used in the massive proof.  Thus the maximal
spacetime factor is the Poincar\'e algebra if the new massless generators are
absent, and the conformal algebra if they are present; in either case it is
accompanied only by an internal symmetry algebra.

\WeinbergSolution{2}
Write \(\partial_\pm=\partial/\partial\sigma^\pm\) and first retain only the
transformation with parameter \(\alpha_2(\sigma^-)\).  Since
\(\partial_+\alpha_2=0\), variation of the bosonic term, followed by the
product rule, gives
\[
\begin{aligned}
 \delta_2\mathcal L_X
 ={}&T\partial_+\!\left(
   \alpha_2\psi_2^\mu\partial_-X_\mu\right)
 +T\partial_-\!\left(
   \alpha_2\psi_2^\mu\partial_+X_\mu\right)\\
 &-2T\alpha_2\psi_2^\mu\partial_+\partial_-X_\mu .
\end{aligned}
\]
The two odd quantities \(\alpha_2\) and \(\psi_2\) anticommute.  Using that
sign when varying the first-order fermion term yields
\[
\begin{aligned}
 \delta_2\mathcal L_{\psi_2}
 &=i(\delta\psi_2^\mu)\partial_+\psi_{2\mu}
   +i\psi_2^\mu\partial_+(\delta\psi_{2\mu})\\
 &=-T\partial_+\!\left(
   \alpha_2\psi_2^\mu\partial_-X_\mu\right)
   +2T\alpha_2\psi_2^\mu\partial_+\partial_-X_\mu .
\end{aligned}
\]
The bulk terms cancel, as do the two \(\partial_+\) boundary terms.  Hence
\[
 \delta_2\mathcal L
 =T\partial_-\!\left(
   \alpha_2\psi_2^\mu\partial_+X_\mu\right).
\]
The \(\alpha_1(\sigma^+)\) transformation is the same calculation with
\(+\leftrightarrow-\) and \(1\leftrightarrow2\):
\[
 \delta_1\mathcal L
 =T\partial_+\!\left(
   \alpha_1\psi_1^\mu\partial_-X_\mu\right).
\]
Therefore the complete change is
\[
 \delta\mathcal L
 =T\partial_+\!\left(\alpha_1\psi_1^\mu\partial_-X_\mu\right)
  +T\partial_-\!\left(\alpha_2\psi_2^\mu\partial_+X_\mu\right).
\]
Its integral vanishes for a closed worldsheet, or for boundary conditions
that make the displayed boundary flux zero.  The Gervais--Sakita action is
therefore invariant without using any field equation.

\WeinbergSolution{3}
First note the sign needed in the auxiliary interaction.  Consistency with
Eq.~\eqref{eq:24.2.11} requires
\[
 gF(A^2-B^2),\qquad
 F=-mA-g(A^2-B^2),
\]
not \(gF(A^2+B^2)\).  This independently verifiable transcription correction
is recorded in \texttt{MODERNIZATION.md} and has been made only in the
exercise edition.

It is economical to organize the direct component calculation with
\[
\begin{gathered}
 P_L=\frac{1+\gamma_5}{2},\qquad
 P_R=\frac{1-\gamma_5}{2},\\
 \phi=\frac{A+iB}{\sqrt2},\qquad
 \mathcal F=\frac{F-iG}{\sqrt2},\qquad
 f(\phi)=\frac m2\phi^2+\frac{\sqrt2g}{3}\phi^3 .
\end{gathered}
\]
Thus \(f'(\phi)=m\phi+\sqrt2g\phi^2\).  In these variables
Eq.~\eqref{eq:24.2.9} is the sum of the kinetic \(D\)-term and the
\(\mathcal F\)-terms of \(f\) and \(f^*\).

Vary every component, integrate terms containing
\(\partial_\mu\delta A\), \(\partial_\mu\delta B\), and
\(\partial_\mu\delta\psi\) once by parts, and use
\(\{\gamma_5,\gamma^\mu\}=0\).  The terms proportional to \(F\), \(G\),
\(m\), and \(g\) without an overall derivative cancel pairwise.  The
three-fermion remainder vanishes by the Majorana interchange identities.
The terms that remain can be displayed without invoking any field equation:
\[
 \delta\mathcal L=\bar\alpha\,\partial_\mu K^\mu ,
\]
where
\[
\begin{aligned}
 K^\mu=\frac1{\sqrt2}\gamma^\mu\Big[
 &-(\gamma^\nu\partial_\nu\phi)\psi_R
  -(\gamma^\nu\partial_\nu\phi^*)\psi_L
  +\mathcal F^*\psi_L+\mathcal F\psi_R\\
 &+2f'(\phi)\psi_L+2f'(\phi)^*\psi_R
 \Big],
\end{aligned}
\]
and \(\psi_{L,R}=P_{L,R}\psi\).  For example, the first line arises from the
kinetic and auxiliary terms, while the second line is the complete mass and
cubic-interaction contribution.  In real fields the same current is
\[
\begin{aligned}
 K^\mu={}&\frac12\gamma^\mu\Big[
   -\gamma^\nu(\partial_\nu A)\psi
   +i\gamma^\nu(\partial_\nu B)\gamma_5\psi
   +F\psi+iG\gamma_5\psi\Big]\\
 &+\gamma^\mu\Big[
   \bigl(mA+g(A^2-B^2)\bigr)\psi
   +i\bigl(mB+2gAB\bigr)\gamma_5\psi\Big].
\end{aligned}
\]
The change of the Lagrangian density is therefore a spacetime divergence.
For fields that vanish sufficiently rapidly at infinity, the boundary
integral is zero and the Wess--Zumino action is supersymmetric off shell.
\end{enumerate}
